\documentclass[10pt,a4paper]{amsart}
\usepackage[margin=19mm]{geometry}
\usepackage{amsmath,amssymb,mathtools}
\usepackage[scr=rsfs,cal=euler]{mathalfa}
\usepackage{xfrac}
\usepackage[unicode,hidelinks,bookmarks=false]{hyperref}
\newcommand{\ie}{i.e.\ }
\newcommand{\cf}{cf.\ }
\newcommand{\resp}{respectively\ }
\newcommand{\Subsection}{\subsection}
\providecommand{\nobreakdash}{\nobreak\hbox{-}}
\setlength{\emergencystretch}{2em}
\multlinegap0pt
\usepackage{mathtools}
\usepackage{amssymb}
\usepackage{bm}
\usepackage{xcolor}
\usepackage[shortlabels]{enumitem}
\newtheorem{proposition}{Proposition}
\title{\fontsize{15.6}{19.5}\selectfont \bf Huber-based Robust System Identification with Near-Optimal Guarantees Across Independent and Adversarial Regimes}
\author{Jihun Kim, Javad Lavaei}
\date{Source: arXiv:2603.27586v2 (CC BY 4.0)}
\begin{document}
\maketitle

\noindent\emph{Source and licence.} \fontsize{15.6}{19.5}\selectfont \bf Huber-based Robust System Identification with Near-Optimal Guarantees Across Independent and Adversarial Regimes. Version arXiv:2603.27586v2 (CC BY 4.0), distributed by arXiv under the Creative Commons Attribution 4.0 International licence, http://creativecommons.org/licenses/by/4.0/, as recorded in the arXiv licence metadata for this exact version and shown on the abstract page https://arxiv.org/abs/2603.27586v2. Full licence notice: \texttt{licenses/arxiv-2603.27586-CC-BY-4.0.txt}.\par\smallskip

\noindent\textit{Editorial scope.} Counted unit: the article's proposition lasso (source0.tex lines 249--255). The article's complete original proof of it is delivered with it (source0.tex lines 257--279, one \texttt{proof} environment of 23 lines). No mathematics is written, repaired or re-derived: the statement and the proof are the article's own lines, in the article's own notation, and no retained line is substituted or re-flowed.  Retained uncounted as the source-local context the proof needs: lines 233--235.  Nothing was omitted from any retained span, so the package carries no omission record.  No retained line contains a citation, so the wrapper carries no bibliography and no bibliography entry.  No retained line contains a figure environment, an image-inclusion command or drawing code, so no image is delivered and the package declares no asset dependency.  Editorial formatting only, in addition: an AMS \texttt{amsart} preamble that restates the article's own declarations and macros used by the retained lines, namely the environments \texttt{proposition} on separate counters, together with the standard packages those lines need.  The retained environments print in this document as Proposition 1 in order of appearance, whereas the article prints the same items with the numbering of its own full document.  The complete native source of the article as distributed in the arXiv e-print for this exact version is bundled unchanged as \texttt{sources/197360/source0.tex}.
\begin{align}\tag{Huber estimator}  \label{hubermin}
    \min_{\{a_i\}_{i=1}^n} \sum_{t=0}^{T-1} \sum_{i=1}^n H_{\mu} (x_{t+1}^{(i)} - a_i^T \phi(x_t)), 
\end{align}

\begin{proposition}
    The problem \eqref{hubermin} is equivalent to the problem 
    \begin{align}\label{lasso}
    \min_{A, \{v_t\}_{t\ge 0}} \sum_{t=0}^{T-1} \frac{1}{2}\|x_{t+1}  - A\phi(x_t)-v_t\|_2^2 + \mu \|v_t\|_1,
\end{align}
where $A$ is a matrix whose rows are $a_i^T$ for $i\in\{1,\dots,n\}$.
\end{proposition}

\begin{proof}
   The joint minimization with respect to $A$ and $\{v_t\}_{t\ge 0}$ is equivalent to  minimizing first with respect to $A$, and then over $\{v_t\}_{t\ge 0}$. Thus, it suffices to show that 
    \begin{equation}\label{eqhuber}
    \begin{split}
       & \min_{\{v_t\}_{t\ge 0}} \sum_{t=0}^{T-1} \frac{1}{2}\|x_{t+1}  - A\phi(x_t)-v_t\|_2^2 + \mu \|v_t\|_1 \\&\hspace{30mm}= \sum_{t=0}^{T-1} \sum_{i=1}^n H_{\mu} (x_{t+1}^{(i)} - a_i^T\phi(x_t))
       \end{split}
    \end{equation}
    for all $A$. Since the left-hand side term is convex in $\{v_t\}_{t\ge 0}$ and can be decoupled along the time $t$ as well as the coordinate $i$, the KKT optimality conditions imply that
    \begin{align*}
       0\in   -(x_{t+1}^{(i)} - a_i^T \phi(x_t)-v_t^{(i)}) + \mu \partial |v_t^{(i)}|, \quad t\ge 0
    \end{align*}
    for every $i\in\{1,\dots,n\}$, where $\partial$ denotes the subderivative.
We consider two cases based on the KKT conditions.

\textbf{Case 1}: $|x_{t+1}^{(i)} - a_i^T \phi(x_t)| \le \mu$. In this case, $v_t^{(i)} = 0$. To see why, note that if $v_t^{(i)} > 0$, the KKT conditions would require $x_{t+1}^{(i)} - a_i^T \phi(x_t) - v_t^{(i)} = \mu$, which incurs a contradiction. The argument for $v_t^{(i)} < 0$ follows similarly.

\textbf{Case 2}: $|x_{t+1}^{(i)} - a_i^T \phi(x_t)| > \mu$. Here, $v_t^{(i)} \ne 0$. If $x_{t+1}^{(i)} - a_i^T \phi(x_t) > \mu$, the KKT conditions imply $v_t^{(i)} = x_{t+1}^{(i)} - a_i^T \phi(x_t) - \mu>0$. Alternatively, if $x_{t+1}^{(i)} - a_i^T \phi(x_t) < -\mu$, the conditions imply $v_t^{(i)} = x_{t+1}^{(i)} - a_i^T \phi(x_t) + \mu<0$.
    
    
    
    
    Substituting the obtained $v_t^{(i)}$ for every $t$ and $i$ back into the left-hand side of \eqref{eqhuber} yields the right-hand side. This completes the proof.
\end{proof}

\end{document}
